\chapter{OffendES y su aplicación en el proyecto}

El corpus OffendES representa un aporte fundamental para la comunidad hispana de PLN, esto es debido a que la mayoría de los recursos disponibles para entrenar y evaluar estos sistemas han sido desarrollados para el idioma inglés, lo que limita significativamente la investigación y el desarrollo de soluciones adaptadas a contextos hispanohablantes.

OffendES, se trata de un conjunto de datos en español diseñado específicamente para la detección de lenguaje ofensivo, compuesto por comentarios recolectados de diferentes redes sociales y anotados manualmente.

\section{Características del corpus OffendES}

\subsection{Fuentes de datos y recolección}

El proceso de recolección de datos se llevó a cabo entre febrero y marzo de 2020, tomando como referencia tres plataformas ampliamente utilizadas por usuarios jóvenes (entre 18 y 24 años): Twitter, Instagram y YouTube. Estas redes sociales fueron seleccionadas por su alto nivel de interacción y su relevancia en el ecosistema digital.

Se eligieron 12 personalidades influyentes del entorno español, seis hombres y seis mujeres, que se caracterizan por su alta exposición mediática y participación en controversias en línea. Para cada uno de ellos, se extrajeron los comentarios correspondientes a sus últimas 50 publicaciones. En el caso de Twitter y YouTube, la recopilación se realizó mediante APIs oficiales; para Instagram, se utilizaron técnicas de web scraping, estableciendo un máximo de 2,000 respuestas por publicación.

Como resultado, se recolectaron un total de 283,622 comentarios, de los cuales 23,788 incluían términos con potencial ofensivo, identificados a partir de lexicones previamente definidos. Posteriormente, se conformó un conjunto final de 60,000 comentarios (23,788 ofensivos y 36,212 no ofensivos), aplicando criterios de diversidad léxica para evitar redundancias y asegurar representatividad.

\subsection{Esquema de anotación}

La anotación del corpus fue realizada mediante la plataforma de \textit{crowdsourcing} Amazon Mechanical Turk (MTurk), empleando únicamente trabajadores ubicados en España para asegurar la comprensión contextual del lenguaje.



\vspace{0.5cm}

Se utilizó un esquema de cinco etiquetas diferenciadas, con el objetivo de capturar distintos matices del lenguaje ofensivo:

\begin{itemize}
    \item \textbf{OFP} (Ofensivo hacia una persona): Comentarios que agreden directamente a un individuo.
    \item \textbf{OFG} (Ofensivo hacia un grupo): Comentarios dirigidos contra colectivos específicos, como género, etnia, orientación sexual o ideología.
    \item \textbf{OFO} (Ofensivo hacia otras entidades): Comentarios que atacan instituciones, marcas, eventos u organizaciones.
    \item \textbf{NOE} (No ofensivo con lenguaje vulgar): Uso de lenguaje soez sin intención de ofender.
    \item \textbf{NO} (No ofensivo): Comentarios neutros o sin lenguaje perjudicial.
\end{itemize}

El corpus se organizó en dos subconjuntos según la cantidad de anotadores por comentario:

\begin{itemize}
    \item \textbf{3-Ann:} Conjunto de 44,951 comentarios, cada uno anotado por tres personas.
    \item \textbf{10-Ann:} Conjunto de 14,989 comentarios, cada uno evaluado por diez anotadores, lo que permite generar vectores de confianza para tareas avanzadas como clasificación multietiqueta.
\end{itemize}

Los atributos presentes en el corpus son los siguientes:

\begin{itemize}
    \item \textbf{comment\_id:} Identificador únic o de cada comentario. \textbf{media:} Plataforma de origen del comentario (Instagram, Twitter o YouTube).
    \item \textbf{influencer:} Nombre del influencer asociado a la publicación.
    \item \textbf{influencer\_gender:} Género del influencer (man o woman). \textbf{comment:} Texto completo del comentario.
    \item \textbf{ann1 -- ann10:} Etiquetas asignadas por cada anotador individual. Los valores posibles son: OFP, OFG, OFO, NOE, NO y None.
    \item \textbf{label:} Etiqueta final asignada a cada comentario, resultante de un proceso de consenso entre los anotadores.
\end{itemize}

\subsection{Distribución y diversidad}

Una característica notable del corpus es su desbalance natural: tanto las etiquetas como las plataformas e influencers no están distribuidos uniformemente. La mayoría de los comentarios provienen de YouTube (alrededor del 75\%), seguido de Instagram (18\%) y Twitter (6\%). Del mismo modo, ciertos influencers como Dalas concentran más del 25\% de los comentarios.

En cuanto a la longitud y variedad del lenguaje, los comentarios muestran diferencias significativas según la plataforma. Por ejemplo, los comentarios de YouTube tienden a ser más extensos (promediode 189 caracteres) y diversos en vocabulario, mientras que los de Instagram son más breves (114 caracteres en promedio) y lexicalmente más pobres. La medida de diversidad léxica también reveló que los comentarios ofensivos (especialmente OFP y NOE) tienden a tener una menor riqueza léxica que los no ofensivos.


